\documentclass{ximera}
\title{On Trails}

\newcommand{\pskip}{\vskip 0.1 in}

\begin{document}
\begin{abstract}
Hiking bugs.
\end{abstract}
\maketitle


\section{Trails}

\begin{question} \label{Qler3r3rfd}
The differentiable function
\[
   h = f(s) \, , \, 0\leq s \leq 5,
\]
expresses the altitude (in km) along a trail in terms of the distance from the trailhead (measured along the trail in km).

\begin{enumerate}
\item Suppose
\[
    \frac{dh}{ds}\Big|_{s=3} = 0.2 .
\]
What are the units of this derivative? Interpret the meaning of the derivative using the language of small changes.

\item What acute angle does the curve $h=f(s)$ make with the $s$-axis at the point $(3, f(3))$.

\item What acute angle does the trail make with the horizontal at the point on the trail $3$ km from the trailhead? 

\item  Suppose
\[
    \frac{dh}{ds}\Big|_{s=4} = 1 .
\]
Explain what this means.
\end{enumerate}

\begin{explanation}
\begin{enumerate}
\item The derivative has units $0.2$ km/km. 

It is a dimensionless scaling factor that converts a small change in the distance to the trailhead from the $3$-km mark to an accurate approximation of the corresponding change in altitude (measured in the same units as the distance from the trailhead, whatever those units may be).

For example, if you walk $1$ meter farther along the trial from the $3$-km mark, your altitude increases by approximately $0.2$ meters.

\item Let $\theta$ be the acute angle (in radians) the curve $h=f(s)$ makes with the $s$-axis at the point $(3, f(3))$. Then 
\[
  \tan \theta = \frac{dh}{ds}\Big|_{s=3} = 0.2 .
\]
And since $-\pi/2 \leq \theta \leq \pi/2$,
\[
 \theta = \arctan(0.2).
\]

\item Let $\phi$ be the acute angle (in radians) the trail makes with the horizontal at the $3$-km mark. From our description in part (a), we know if you walk $\Delta s\sim 0$ meters farther \emph{along the trail} your altitude increases by approximately $0.2 \Delta h$ meters. This tells us
\[
   \sin\phi \sim \frac{\Delta h}{\Delta s} \sim 0.2 .
\]
Taking a limit as $\Delta s \to 0$ shows that
\[
   \sin \phi = \frac{dh}{ds}\Big|_{s=3} = 0.2 .
\]
Since $-\pi/2 \leq \phi \leq \pi/2$, at the $3$-km mark the trail slopes upward at an angle of
\[
    \phi = \arcsin(0.2) \text{ rad}
\]
relative to the horizontal.

\item This means that the trial is vertical at the $4$-km mark. You would be rock-climbing at this point of the trail, not hiking.
\end{enumerate}
\end{explanation}

\end{question}


\begin{question} \label{Wer3fds}
The function
\[
    h = f\left(s\right)=2000+\frac{1800}{\pi}\tan\left(\frac{\pi}{12}s\right), 0\leq s\leq 4.2,
\]
expresses the altitude of a trail (in meters) in terms of the distance from the trailhead (in km).
\begin{enumerate}
\item What angle does the trail make with the horizontal at the $4$-km mark?
\item At what altitude does the trail slope upward at a $30^\circ$ angle relative to the horizontal?
\end{enumerate}
\end{question}




\begin{question} \label{Qdf3rra}
The function 
\[
 h = f(s) \, , \, 0\leq s \leq 15,
\]
expresses the alitude (in meters) of a beetle's path in terms of the distance along the path (in meters) from the start. Its graph is shown below (blue).

The profile of the actual trail (red) is also graphed below.

Point $P$ with coordinates $(s,h) = (s,f(s))$ lies on the graph of the altitude function. Point $B$ is the corrresponding point on the trail. It is $s$ km from the trailhead and has altitude $h=f(s)$.

\begin{enumerate}

\item Use the graphs to explain the following:
\begin{enumerate}
\item Why and how you would use the inverse tangent function to compute the acute angle the curve $h=f(s)$ makes with the $s$-axis at $P$.

\item Why and how you would use the inverse sine function to compute the acute angle the \emph{trail} makes with the horizontal at $B$.
\end{enumerate}

\begin{freeResponse}
\end{freeResponse}

\item Suppose 
\[
 h = f(s) = 5 \sin \left( \frac{s}{5} \right) \, , \, 0\leq s\leq 15.
\]

\begin{enumerate}
\item Find a function 
\[
   \theta = g(s) \, , \, 0\leq s\leq 15
\]
that expresses the acute angle the trail makes with the horizontal in terms of the distance from the start of the trail. 

\item At what point on the trail does the trail slope downward at an angle of $\arcsin(4/5)$ below the horizontal?

\item Evaluate the derivative $d\theta/ds$ at the point on the trail where it slopes downward at an angle of $\arcsin(4/5)$ below the horizontal. Include units.

Then interpret the meaning of the derivative as a scaling factor in terms of two \emph{specific} small changes.


\end{enumerate}
\end{enumerate}
\begin{onlineOnly}
    \begin{center}
\desmos{lktn53d7ju}{900}{600}
\end{center}
\end{onlineOnly}

\href{https://www.desmos.com/calculator/lktn53d7ju}{151: Hiking Trail 54C}



\end{question}

\section{Trails Described with Functions}

In this section we model a trail with a curve $h=f(x)$, where $x$ and $h$ are measured in meters. We take the $h$-coordinate to measure the height (in meters) above sea-level and suppose the positive $x$-axis points due east.

\begin{question}  \label{QOideIIEIEII}
A beetle crawls eastward along the trail $h = f(x)$ at a constant speed of $v_0$ meters/sec.

At what rate is the beetle gaining altitude if

\begin{enumerate}
\item $f(x) = x$ ?

\item $f(x)=2x$ ?

\item $f(x) = mx$ ? 

\item Find the limit of the beetle's rate of ascent as $m\to \infty$ as it crawls on the trail $h=mx$ at the constant speed of $v_0$ meters/sec.
\end{enumerate}
 
\end{question}

\begin{question}  \label{QPeriERE}
\begin{enumerate}

\item A beetle crawls along the parabola
\[
     h = 9-\frac{1}{4}\left(x-6\right)^{2} \, , \, 0\leq x \leq 12,
\]
(coordinates in meters) in the eastward direction at a constant speed of $v_0$ m/sec. At what rate is it ascending when it is $10$ meters above the ground for the second time?

\item A beetle crawls along the curve
\[
     h =4-3\cos\left(\frac{x}{2}\right) \, , \, 0\leq x \leq 12,
\]
(coordinates in meters) in the eastward direction at a constant speed of $v_0$ m/sec. At what rate is it ascending when it is $5$ meters above the ground for the second time?

\end{enumerate}

\end{question}




\section{Parametrically Defined Trails}


\end{document}